% 02-cesm3-cam7 — "CESM3/CAM7". \section = \frametitle: each titled slide gets its own
% section so the Berkeley sidebar lists the deck by slide title.
% Science-level configuration only (no code identifiers). Facts: SPEC.md,
% docs/monthly-2008-case-plan.md, RUN.md "Met nudging".
\section{CESM3/CAM7}
\begin{frame}[t]
\frametitle{CESM3/CAM7}
\footnotesize
\begin{columns}[T]
\begin{column}{0.49\textwidth}
\begin{block}{Model}
\begin{itemize}
  \item CESM3 beta08, CAM7 physics,\\RRTMGP radiation
  \item MAM4 modal aerosols (four modes)
  \item Compset F2000dev: year-2000 forcing,\\prescribed SST and sea ice
  \item f09 ($0.9^\circ \times 1.25^\circ$), 32 levels
\end{itemize}
\end{block}
\begin{block}{Constraints}
\begin{itemize}
  \item Meteorology: hourly ERA5 winds,\\temperature and humidity
  \item Aerosol: 3-hourly MODIS AOD\\at 550 nm
\end{itemize}
\end{block}
\end{column}
\begin{column}{0.49\textwidth}
\begin{block}{2008 Test Run}
\begin{itemize}
  \item Twelve monthly cases,\\10-day spin-up each
  \item 3-hourly history and band optics
  \item Derecho, one node per case
\end{itemize}
\end{block}
\begin{block}{Additions to CAM7}
\begin{itemize}
  \item AOD constraint applied\\before each radiation call
  \item 3-D aerosol optics written\\in the 26 SARB bands
  \item Off by default;\\otherwise unmodified CAM7
\end{itemize}
\end{block}
\end{column}
\end{columns}
\end{frame}
